\clearpage
\subsection{Page-Table Entry Format}\label{section:page-table-entry-format}

Every page-table entry is one naturally aligned 64-bit descriptor. Bits
\texttt{55..12} are a fixed physical frame number field, so this format has a
56-bit architectural physical-address ceiling. An implementation reports its
exact \texttt{PABITS} value, from 32 through 56, through CPUID
\texttt{ADDRESS\_WIDTHS}. In a present descriptor, PFN bits at or above that
value must be zero. Bits \texttt{57..56} are \texttt{SW[1:0]} and are ignored
by hardware; bits \texttt{63..58} are reserved and must be zero.

\texttt{P} in bit 0 is common. When \texttt{P=0}, the remaining bits are
ignored and the entry is not present. When \texttt{P=1}, \texttt{T} in bit 1
selects the leaf or table-pointer layout. A nonzero reserved bit in a present
descriptor raises \texttt{PAGE\_FAULT.MALFORMED\_PTE}.

\manualtablecaption{Leaf Descriptor Fields (P=1, T=0)}
\begin{manualtabular}{@{}p{0.75in}p{0.75in}p{4.00in}@{}}
\toprule
\rowcolor{ManualHeaderFill}\textbf{Field} & \textbf{Bits} & \textbf{Meaning}\\
\midrule
\texttt{P} & \texttt{0} & present, one\\
\texttt{T} & \texttt{1} & leaf selector, zero\\
\texttt{AM} & \texttt{4..2} & permission and Normal/MMIO access class\\
\texttt{U} & \texttt{5} & user-domain permission\\
\texttt{G} & \texttt{6} & global translation\\
\texttt{A} & \texttt{7} & accessed state\\
\texttt{D} & \texttt{8} & exact dirty state\\
\texttt{CP} & \texttt{10..9} & cache policy\\
reserved & \texttt{11} & must be zero\\
\texttt{PFN} & \texttt{55..12} & mapped physical frame number\\
\texttt{SW[1:0]} & \texttt{57..56} & software-defined, ignored by hardware\\
reserved & \texttt{63..58} & must be zero\\
\bottomrule
\end{manualtabular}

\manualtablecaption{Leaf AM Encodings}
\begin{manualtabular}{@{}p{0.65in}p{1.35in}p{1.15in}p{2.35in}@{}}
\toprule
\rowcolor{ManualHeaderFill}\textbf{AM} & \textbf{Permission} & \textbf{Class} & \textbf{Name}\\
\midrule
\texttt{000} & read & Normal & \texttt{R}\\
\texttt{001} & write & Normal & \texttt{W}\\
\texttt{010} & execute & Normal & \texttt{X}\\
\texttt{011} & read/write & Normal & \texttt{RW}\\
\texttt{100} & read/execute & Normal & \texttt{RX}\\
\texttt{101} & read & MMIO & \texttt{MMIO\_R}\\
\texttt{110} & write & MMIO & \texttt{MMIO\_W}\\
\texttt{111} & read/write & MMIO & \texttt{MMIO\_RW}\\
\bottomrule
\end{manualtabular}

No leaf encoding grants both write and execute, and no MMIO leaf grants
execute. This W\textasciicircum X property applies to one selected translation.
The architecture does not compare aliases: distinct descriptors for the same
physical page may have different permissions, AM classes, or CP values.

\manualtablecaption{Table-Pointer Descriptor Fields (P=1, T=1)}
\begin{manualtabular}{@{}p{0.75in}p{0.75in}p{4.00in}@{}}
\toprule
\rowcolor{ManualHeaderFill}\textbf{Field} & \textbf{Bits} & \textbf{Meaning}\\
\midrule
\texttt{P} & \texttt{0} & present, one\\
\texttt{T} & \texttt{1} & table-pointer selector, one\\
\texttt{R/W/X} & \texttt{2/3/4} & subtree maximum permissions\\
\texttt{U} & \texttt{5} & subtree maximum user-domain permission\\
reserved & \texttt{6} & must be zero\\
\texttt{A} & \texttt{7} & accessed state\\
reserved & \texttt{11..8} & must be zero\\
\texttt{PFN} & \texttt{55..12} & next-table physical frame number\\
\texttt{SW[1:0]} & \texttt{57..56} & software-defined, ignored by hardware\\
reserved & \texttt{63..58} & must be zero\\
\bottomrule
\end{manualtabular}

A table pointer with \texttt{R=W=X=0}, or a table pointer at L1, is malformed.
The walker begins with R, W, X, and U enabled, intersects each table bitmap and
U bit, then intersects the terminating leaf\textquotesingle{}s decoded AM permissions and U bit.
The leaf alone supplies G, CP, and Normal/MMIO class. A leaf at level $L$ must
have a physical base aligned to $2^{12+9(L-1)}$ bytes.

The CP field is orthogonal to AM: all 32 AM/CP combinations are structurally
valid. CP values 0 through 3 mean cacheable write-back, coherent uncacheable,
coherent write-through, and coherent write-combining respectively. CP never
weakens an MMIO execution rule.

Hardware sets A on each structurally valid descriptor consumed by a completed
walk and may also set A during a speculative walk whose instruction later
faults or is squashed. D is set only when the corresponding store or atomic update
commits, and every required D update succeeds before its memory update becomes
visible. Hardware A/D changes are conditional atomic updates of the complete
64-bit PTE and preserve every other field. A concurrent change causes a
re-read and complete revalidation; stale fields are never written back.
\texttt{PTQUERY} and \texttt{VTOP} do not modify A or D.
